
\documentclass{article}
\usepackage{colortbl}
\usepackage{makecell}
\usepackage{multirow}
\usepackage{supertabular}

\begin{document}

\newcounter{utterance}

\centering \large Interaction Transcript for game `codenames', experiment `opponent\_difficult', episode 4 with reflexion{-}Qwen3.5{-}27B.
\vspace{24pt}

{ \footnotesize  \setcounter{utterance}{1}
\setlength{\tabcolsep}{0pt}
\begin{supertabular}{c@{$\;$}|p{.15\linewidth}@{}p{.15\linewidth}p{.15\linewidth}p{.15\linewidth}p{.15\linewidth}p{.15\linewidth}}
    \# & $\;$A & \multicolumn{4}{c}{Game Master} & $\;\:$B\\
    \hline

    \theutterance \stepcounter{utterance}  
    & & & \multicolumn{2}{p{0.3\linewidth}}{
        \cellcolor[rgb]{0.9,0.9,0.9}{
            \makecell[{{p{\linewidth}}}]{
                \texttt{\tiny{[GM$|$GM]}}
                {'hidden': {'team': ['robin', 'unicorn', 'smuggler', 'grace'], 'innocent': ['figure', 'drop', 'queen', 'bank', 'scorpion', 'hood'], 'opponent': [], 'assassin': ['net']}, 'revealed': {'team': {'team': [], 'innocent': [], 'opponent': [], 'assassin': []}, 'opponent': {'team': [], 'innocent': [], 'opponent': [], 'assassin': []}}}
            }
        }
    }
    & & \\ \\

    \theutterance \stepcounter{utterance}  
    & & \multicolumn{4}{p{0.6\linewidth}}{
        \cellcolor[rgb]{0.9,0.9,0.9}{
            \makecell[{{p{\linewidth}}}]{
                \texttt{\tiny{[P1$\langle$GM]}}
                \texttt{You have attempted this same task before and did not succeed. These are the plans you wrote after those attempts.} \\
\texttt{Trial \#1: New plan: My previous strategy failed because I prioritized clearing words in small, isolated groups (like "mercury" alone or "robin" alone), which allowed the opponent to catch up and eventually win. Specifically, I missed the opportunity to group "unicorn" and "grace" with "robin" under a broader clue like "myth" or "fable" earlier, and I failed to link "smuggler" with "robin" (both are birds/characters in specific contexts) or "grace" with "mercury" (both are names of planets/gods) to maximize turns. Next time, I will analyze the full list of remaining words at the start of each turn to find the single clue that connects the maximum number of targets, even if the association is slightly more abstract, to clear multiple words in one turn and outpace the opponent. I will also avoid clues that are too specific to a single word unless it is the very last word remaining.} \\
\\ 
\texttt{Trial \#2: New plan: My previous attempt failed because I selected a clue ("MYTH") that was too broad and ambiguous, causing my partner to likely guess a distractor (like "figure" or "queen") or miss the connection entirely, ending our turn prematurely. Specifically, "robin" is a real bird, not a myth, making the association weak and risky. Next time, I will prioritize high{-}confidence, direct semantic links over broad categories. I will look for a clue that strongly connects at least two words without risking a distractor. For this specific set, I will use "BIRD" to target "robin" and "smuggler" (referencing the bird "smuggler" or a specific character association if valid, but more safely, I will split the turn: first use "BIRD" for "robin" and "smuggler" if the context allows, or better yet, use "HORSE" for "unicorn" and "grace" (as in Grace Kelly's horse or a specific reference) or simply "FABLE" for "unicorn" and "robin" (Pied Piper/Robin Hood context) while strictly avoiding "MYTH" which includes non{-}mythical birds. Actually, the safest high{-}yield move is to use "BIRD" for "robin" and "smuggler" (if "smuggler" is a bird type in this context, otherwise just "robin") and "HORSE" for "unicorn". Wait, "smuggler" is not a bird. I will use "BIRD" for "robin" only if no other link exists, but better: "FABLE" for "unicorn" and "robin" (Robin Hood is a legend, Unicorn is a legend). No, "robin" is a bird. The best link is "LEGEND" for "unicorn" and "robin" (Robin Hood). But "grace" and "smuggler" are left. I will instead use "NAME" for "grace" and "robin" (common names) and "smuggler" (maybe not). Let's try "ANIMAL" for "robin" and "unicorn". That is safe. "ANIMAL" connects "robin" and "unicorn". "smuggler" and "grace" are harder. I will use "ANIMAL" to get "robin" and "unicorn" safely, avoiding the "MYTH" trap. If that fails, I will use "BIRD" for "robin" and "HORSE" for "unicorn" in separate turns to ensure no distractor guesses. The key is to avoid clues that force the partner to guess a distractor like "figure" or "queen". I will choose "ANIMAL" for "robin, unicorn" as the primary strategy.} \\
\\ 
\\ 
\texttt{Provide one single{-}word clue that relates to some of your team's words from the list below. You can choose to how many and to which words your clue relates to.} \\
\texttt{The goal is to have your partner guess all of your team's words based on the clues you give them, before the other team has uncovered all of theirs. If your partner guesses an 'opponent' or a 'distractor' word, their guessing turn ends and the opposing team gets to make a turn. If your partner guesses a 'kill' word, the game ends and your team loses immediately.} \\
\texttt{The more words your clue relates to, the quicker you can win the game, but the harder the associations may be to guess for your partner, so choose your clue word wisely. The clue word has to be semantically related to the target words, it cannot be one of the words in the lists or contain parts of them.} \\
\\ 
\texttt{Always give your single{-}word clue and your comma{-}separated list of related target words in the following format and make your answers as short as possible, never include any other text than is required in this form:} \\
\\ 
\texttt{CLUE: <WORD>} \\
\texttt{TARGETS: <TARGETS>} \\
\\ 
\texttt{Your team words are: robin, unicorn, smuggler, grace.} \\
\texttt{Your opponent's team words are: .} \\
\texttt{Distractor words are: figure, drop, queen, bank, scorpion, hood.} \\
\texttt{The kill words are: net.} \\
            }
        }
    }
    & \\ \\

    \theutterance \stepcounter{utterance}  
    & \multicolumn{4}{p{0.6\linewidth}}{
        \cellcolor[rgb]{0.9,0.9,0.9}{
            \makecell[{{p{\linewidth}}}]{
                \texttt{\tiny{[P1$\rangle$GM]}}
                \texttt{CLUE: NAME} \\
\texttt{TARGETS: robin, grace} \\
            }
        }
    }
    & & \\ \\

    \theutterance \stepcounter{utterance}  
    & & & \multicolumn{2}{p{0.3\linewidth}}{
        \cellcolor[rgb]{0.9,0.9,0.9}{
            \makecell[{{p{\linewidth}}}]{
                \texttt{\tiny{[GM$|$GM]}}
                \texttt{name} \\
            }
        }
    }
    & & \\ \\

    \theutterance \stepcounter{utterance}  
    & & & \multicolumn{2}{p{0.3\linewidth}}{
        \cellcolor[rgb]{0.9,0.9,0.9}{
            \makecell[{{p{\linewidth}}}]{
                \texttt{\tiny{[GM$|$GM]}}
                ['robin', 'grace']
            }
        }
    }
    & & \\ \\

    \theutterance \stepcounter{utterance}  
    & & & \multicolumn{2}{p{0.3\linewidth}}{
        \cellcolor[rgb]{0.9,0.9,0.9}{
            \makecell[{{p{\linewidth}}}]{
                \texttt{\tiny{[GM$|$GM]}}
                {'word': 'robin', 'assignment': 'team'}
            }
        }
    }
    & & \\ \\

    \theutterance \stepcounter{utterance}  
    & & & \multicolumn{2}{p{0.3\linewidth}}{
        \cellcolor[rgb]{0.9,0.9,0.9}{
            \makecell[{{p{\linewidth}}}]{
                \texttt{\tiny{[GM$|$GM]}}
                {'word': 'grace', 'assignment': 'team'}
            }
        }
    }
    & & \\ \\

    \theutterance \stepcounter{utterance}  
    & & & \multicolumn{2}{p{0.3\linewidth}}{
        \cellcolor[rgb]{0.9,0.9,0.9}{
            \makecell[{{p{\linewidth}}}]{
                \texttt{\tiny{[GM$|$GM]}}
                \texttt{opponent has won} \\
            }
        }
    }
    & & \\ \\

    \theutterance \stepcounter{utterance}  
    & \multicolumn{4}{p{0.6\linewidth}}{
        \cellcolor[rgb]{0.95,0.95,0.95}{
            \makecell[{{p{\linewidth}}}]{
                \texttt{\tiny{[P2$\rightarrow$P1]}}
                \texttt{Post{-}episode Reflexion (not sent to GM):} \\
\texttt{New plan: My previous attempt failed because I chose a clue ("NAME") that was too broad and ambiguous, likely causing my partner to guess a distractor (like "queen" or "figure") or miss the connection entirely, ending our turn prematurely. Specifically, "robin" and "grace" are common names, but "smuggler" and "unicorn" are not, leaving half the board untouched while risking a turn loss. Next time, I will prioritize high{-}confidence, direct semantic links that connect the maximum number of words without ambiguity. For this specific set, I will use "ANIMAL" to safely target "robin" and "unicorn" (both are creatures), or "LEGEND" to target "unicorn" and "robin" (referencing Robin Hood and the mythical unicorn), ensuring the clue is distinct from distractors like "queen" or "figure". I will avoid clues that rely on weak associations like "NAME" which apply to only a subset of words and invite risky guesses.} \\
            }
        }
    }
    & & \\ \\

\end{supertabular}
}

\end{document}
